\documentclass[
    left=2.5cm,
    right=2.5cm,
    top=2.5cm,
    bottom=3cm,
    bindingoffset=6mm,
    nohyphenation=false
]{szablon/eiti/sprawozdanie}

\langpol
\graphicspath{{rysunki/}}
\addbibresource{bibliografia.bib}

% Dodanie usuwa jedno z ostrzeżeń
\usepackage{csquotes}
\usepackage{xurl}
% Dodanie tego wpływa na wyświetlanie zdjęć w odpowiedni sposób przy adnotacji [H]
\usepackage{float}
\usepackage{subfig}
\usepackage{amsmath}
\usepackage[most]{tcolorbox}
\usepackage{xurl}
\usepackage{hyperref}
\usepackage{booktabs}

\newtcolorbox{githubbox}{
  colback=gray!5,
  colframe=black!35,
  arc=2mm,
  boxrule=0.4pt,
  left=2mm,
  right=2mm,
  top=1mm,
  bottom=1mm
}

\begin{document}

\input{rozdzialy/0-strona-tytulowa}

\newpage
\pagestyle{headings}

\newpage
\section{Opis rozwiązania}

Celem projektu było ulepszenie procesu inwersji DDIM\cite{song2021denoising} dla modelu Stable
Diffusion XL \cite{podell2024sdxl}. W klasycznym użyciu model dyfuzyjny zaczyna od losowego
szumu i wykonuje kolejne kroki odszumiania, aż do otrzymania latentu
odpowiadającego obrazowi. Inwersja wykonuje operację przeciwną: dla danego
obrazu albo latentu obrazu próbuje odtworzyć taki latent szumowy, z którego
ponowne próbkowanie wygeneruje możliwie podobny obraz. Problem polega na
tym, że standardowa inwersja DDIM korzysta z przybliżenia lokalnego. Dla
przejścia pomiędzy kolejnymi stanami trajektorii zakłada ona, że predykcja
szumu obliczona dla czystszego latentu może zastąpić predykcję, którą model
otrzymałby dla następnego, bardziej zaszumionego stanu. To przybliżenie jest
wygodne obliczeniowo, ale powoduje błąd akumulujący się w kolejnych krokach
inwersji.

Zaimplementowane rozwiązanie uczy adapter LoRA\cite{hu2022lora} dla UNetu modelu SDXL tak,
aby poprawić właśnie tę predykcję wykorzystywaną w inwersji. Bazowy model
SDXL pozostaje zamrożony i pełni rolę nauczyciela, natomiast trenowane są
wyłącznie niskowymiarowe macierze LoRA dodane do warstw uwagi UNetu. Celem
uczenia nie jest bezpośrednie przewidywanie kolejnego latentu trajektorii,
ale przewidywanie szumu nauczyciela dla następnego, bardziej zaszumionego
stanu:
\[
    \varepsilon_{\phi}(x_{t-1}, t, c)
    \approx
    \varepsilon_{\theta}(x_t, t, c),
\]
gdzie $\theta$ oznacza zamrożony bazowy UNet SDXL, $\phi$ oznacza UNet z
trenowanym adapterem LoRA, $x_{t-1}$ jest czystszym latentem dostępnym w
trakcie inwersji, $x_t$ jest odpowiadającym mu bardziej zaszumionym latentem,
a $c$ oznacza warunkowanie tekstowe SDXL.

\subsection{Przygotowanie danych}

Dane treningowe są budowane z promptów pochodzących ze zbioru Recap-COCO\cite{li2025what}.
Dla każdego promptu uruchamiane jest próbkowanie modelem 
\url{stabilityai/stable-diffusion-xl-base-1.0} z harmonogramem DDIM.
Konfiguracja projektu używa obrazów $1024 \times 1024$, 50 kroków
inferencji oraz skali guidance\cite{DBLP:journals/corr/abs-2207-12598} równej 1. Dla każdej próbki zapisywany jest
obraz końcowy, treść promptu, metadane, lista timestepów oraz pełna
trajektoria latentów w pliku \texttt{trajectory.pt}. Trajektoria zapisana
jest w kolejności próbkowania: od początkowego szumu do coraz czystszych
latentów prowadzących do obrazu.

Oprócz samej trajektorii zapisywane są tensory potrzebne do treningu
adaptera. Dla SDXL warunkowanie nie składa się wyłącznie z
\url{prompt\_embeds} przekazywanego jako \url{encoder\_hidden\_states}.
Model wymaga również dodatkowego warunkowania:
\url{pooled\_prompt\_embeds} oraz \url{add\_time\_ids}, przekazywanych
do UNetu przez \url{added\_cond\_kwargs}. Implementacja zapisuje więc dla
każdej próbki plik \url{conditioning.pt} zawierający:
\url{prompt\_embeds}, \url{pooled\_prompt\_embeds} i
\url{add\_time\_ids}. Dzięki temu podczas treningu nie trzeba ponownie
uruchamiać enkoderów tekstu.

Drugim zapisanym elementem jest \texttt{target\_eps.pt}. Jest to stos
predykcji szumu obliczonych przez zamrożony bazowy UNet SDXL dla kolejnych
bardziej zaszumionych latentów trajektorii. Dla trajektorii zapisanej jako
\[
    z_0, z_1, \ldots, z_N,
\]
gdzie $z_0$ jest początkowym szumem, a $z_N$ najczystszym latentem obrazu,
próbka treningowa o indeksie $i$ używa $z_{i+1}$ jako wejścia ucznia i
predykcji nauczyciela dla $z_i$ jako celu. Ostatni latent trajektorii nie
jest używany jako osobne wejście, ponieważ nie ma po nim kolejnego celu.

\subsection{Zbiór treningowy i funkcja straty}

Klasa \texttt{LatentTrajectoryDataset} traktuje każdą zapisaną trajektorię
jako zbiór przejść pomiędzy kolejnymi latentami. Dla pojedynczego elementu
zwracane są: czystszy latent \texttt{x\_clean}, timestep celu,
\texttt{prompt\_embeds}, \texttt{pooled\_prompt\_embeds},
\texttt{add\_time\_ids} oraz zapisany wcześniej \texttt{target\_eps}. Dzięki
temu trening nie musi obliczać celu nauczyciela w każdej iteracji, co
zmniejsza koszt pojedynczego kroku.

Podczas treningu UNet z aktywnym adapterem LoRA dostaje czystszy latent oraz
ten sam timestep i to samo warunkowanie, które odpowiadają celowi
nauczyciela. Strata jest średnim błędem kwadratowym pomiędzy predykcją
ucznia a zapisanym celem:
\[
    \mathcal{L}(\phi)
    =
    \left\|
        \varepsilon_{\phi}(z_{i+1}, t_i, c)
        -
        \varepsilon_{\theta}(z_i, t_i, c)
    \right\|_2^2.
\]
Jest to bezpośrednie dopasowanie do przybliżenia wykorzystywanego przez
DDIM inversion: model ma lepiej estymować predykcję dla bardziej
zaszumionego stanu, mając dostęp tylko do czystszego latentu.

\subsection{Trening adaptera LoRA}

Trening został zaimplementowany w klasie \texttt{SDXLInversionTrainer}.
Pipeline SDXL jest ładowany w trybie \texttt{fp16}/\texttt{bf16}, a jego
składniki są zamrażane. Adapter LoRA jest wstrzykiwany wyłącznie do UNetu,
w warstwach uwagi \texttt{to\_q}, \texttt{to\_k}, \texttt{to\_v} oraz
\texttt{to\_out.0}. Trenowane są tylko parametry LoRA, natomiast bazowe wagi
modelu, VAE oraz enkodery tekstu pozostają niezmienione. Takie podejście
ogranicza liczbę trenowanych parametrów i pozwala używać adaptera tylko w
inwersji, bez wpływania na standardowe generowanie obrazów, jeśli adapter
nie zostanie jawnie aktywowany.

W projekcie przygotowano trzy warianty adaptera różniące się rzędem macierzy
LoRA: $r=8$, $r=16$ oraz $r=32$. Dla nich ustawiono odpowiednio
\texttt{lora\_alpha} równe 4, 8 i 16. Wspólna konfiguracja treningu używa
optymalizatora AdamW\cite{DBLP:journals/corr/abs-1711-05101}, współczynnika uczenia $5 \cdot 10^{-5}$, zerowego
weight decay, gradient clipping do wartości 1 oraz harmonogramu cosine z
1000 krokami rozgrzewki. Checkpointy zapisywane są cyklicznie jako same wagi adaptera
oraz, opcjonalnie, pełny stan treningu zawierający stan optymalizatora,
schedulera i generatorów losowych.

\subsection{Użycie adaptera w inwersji i ewaluacja}

Wytrenowany adapter jest używany podczas etapu inwersji. Skrypt ewaluacyjny
ładuje zapisany latent końcowy obrazu, koduje prompt, aktywuje LoRA w UNecie
i wykonuje DDIM inverse scheduler, zapisując odtworzony latent szumowy
\texttt{inverted\_noise.pt}. Następnie osobny etap rekonstrukcji startuje z
tego latentu szumowego i wykonuje zwykłe próbkowanie DDIM do obrazu
\texttt{reconstructed.png}. W konfiguracji ewaluacji rekonstrukcja może być
uruchomiona z wyłączonym adapterem LoRA, aby odizolować wpływ adaptera na
samą inwersję, a nie na standardowy proces generowania.

Raport ewaluacyjny porównuje odtworzony szum z oryginalnym początkowym
szumem trajektorii oraz obraz zrekonstruowany z obrazem końcowym zapisanym
podczas generowania danych. Liczone są między innymi metryki MSE, RMSE, MAE,
PSNR i SSIM dla obrazów, statystyki błędu latentów, podobieństwo
odtworzonego szumu do rozkładu normalnego oraz wizualizacje błędów obrazu.
Taka ewaluacja sprawdza więc dwa aspekty rozwiązania: czy LoRA poprawia
odtworzenie latentnego szumu oraz czy po ponownym próbkowaniu zachowany jest
obraz, od którego zaczęła się inwersja.
\begin{githubbox}
\textbf{Kod źródłowy projektu:}\\
\href{https://github.com/AtarSander/Diffusion-Inversion-for-Image-Editing}{\nolinkurl{https://github.com/AtarSander/Diffusion-Inversion-for-Image-Editing}}
\end{githubbox}

\section{Wyniki eksperymentów}

Celem eksperymentów było sprawdzenie, czy adapter LoRA uczony do korekty
predykcji szumu może poprawić jakość inwersji DDIM. Ocena została
przeprowadzona w trzech aspektach: podobieństwa odtworzonego szumu do
rzeczywistego szumu początkowego, jakości rekonstrukcji obrazu oraz
przydatności uzyskanych latentów w edycji tekstowej. Szczególnie istotne
było sprawdzenie, czy poprawa własności latentów nie odbywa się kosztem
rekonstrukcji obrazu.

Zbiór ewaluacyjny obejmował 1522 próbki walidacyjne ze zbioru recap-coco przygotowane analogicznie
do danych treningowych.

Porównanie wykonano względem standardowej DDIM inversion. Dodatkowo
uwzględniono warianty, w których adapter LoRA był aktywny tylko przez
pierwsze kroki inwersji, ponieważ zgodnie z założeniami projektu największy
błąd aproksymacji powinien pojawiać się na początku trajektorii.

\subsection{Porównywane warianty}

Porównano następujące konfiguracje inwersji:
\begin{itemize}
    \item \textbf{DDIM Inv.} -- standardowa inwersja DDIM bez adaptera LoRA.
    \item \textbf{LoRA 1--5 kroków} -- adapter LoRA aktywny wyłącznie przez
          pierwsze $1,\ldots,5$ kroków inwersji, czyli odpowiednio
          $2\%,\ldots,10\%$ trajektorii.
    \item \textbf{Pełna LoRA} -- adapter aktywny przez wszystkie 50 kroków
          inwersji.
\end{itemize}

Adapter był aktywny tylko podczas inwersji. Rekonstrukcja obrazu oraz
generowanie obrazu edytowanego były wykonywane bez LoRA. Dzięki temu
porównanie mierzy jakość otrzymanego latentu, a nie zmianę samego procesu
generacji.


\subsection{Normalność latentów i jakość rekonstrukcji}

Pierwszym celem było sprawdzenie, czy adapter zmniejsza odchylenie
odtworzonego szumu od rzeczywistego szumu początkowego. Do opisu własności
latentów wykorzystano średnią korelację Pearsona w patchach $8 \times 8$.
Im niższa wartość tej korelacji, tym słabsze lokalne zależności przestrzenne
w latencie, a więc tym bliżej zachowuje się on jak niezależnie próbkowany
szum gaussowski. Drugą metryką normalności jest dywergencja KL liczona
niezależnie dla każdej pozycji latentu i uśredniana po wymiarach.

\subsubsection{Pełna aktywacja LoRA}

Pierwszym porównaniem jest pełna aktywacja adaptera LoRA dla trzech
wariantów rzędu: $r=8$, $r=16$ oraz $r=32$ z inwersją DDIM.

\begin{table}[H]
    \centering
    \caption{Porównanie adapterów LoRA dla różnych rzędów z inwersją DDIM.}
    \label{tab:lora-rank-comparison}
    \resizebox{\textwidth}{!}{%
    \begin{tabular}{lrrrrrr}
        \toprule
        Metoda &
        Corr. $\downarrow$ &
        KL Div. $\times 10^2$ $\downarrow$ &
        Recon. MAE $\downarrow$ &
        LPIPS $\downarrow$ &
        PSNR $\uparrow$ &
        SSIM $\uparrow$ \\
        \midrule
        DDIM Inv. & 0.1601 & 0.8566 & 0.02663 & 0.09617 & 28.32 & 0.8112 \\
        \midrule
        $r=8$, $\alpha=4$   & 0.0999 & 0.0767 & 0.02063 & 0.05982 & 30.06 & 0.8654 \\
        $r=16$, $\alpha=8$  & \underline{0.0989} & \underline{0.0671} & \underline{0.02028} & \underline{0.05774} & \underline{30.22} & \underline{0.8699} \\
        $r=32$, $\alpha=16$ & \textbf{0.0972} & \textbf{0.0504} & \textbf{0.01970} &  \textbf{0.05538} & \textbf{30.40} & \textbf{0.8743} \\
        \bottomrule
    \end{tabular}}
\end{table}

Najlepsze wyniki dla pełnej aktywacji uzyskał adapter $r=32$. W porównaniu
z $r=16$ zmniejsza on korelację z $0.0989$ do $0.0972$, KL Div.
$\times 10^2$ z $0.0671$ do $0.0504$. Jednocześnie poprawia metryki rekonstrukcji: Recon. MAE spada do
$0.01970$, LPIPS\cite{DBLP:journals/corr/abs-1801-03924} do $0.05538$, PSNR rośnie do $30.40$, a SSIM do $0.8743$.
Wariant $r=16$ jest drugim najlepszym ustawieniem w większości metryk,
natomiast $r=8$ daje najsłabszą, ale nadal wyraźną poprawę względem DDIM
Inv.


Rysunek~\ref{fig:qualitative-reconstruction} przedstawia przykładowe wyniki
rekonstrukcji. W każdej grupie lewa kolumna pokazuje obraz źródłowy, środkowa
kolumna rekonstrukcję po standardowej inwersji DDIM, a prawa kolumna
rekonstrukcję po inwersji z adapterem LoRA $r=32$.

\begin{figure}[H]
    \centering
    \subfloat{%
        \includegraphics[width=0.31\textwidth]{images/1final.png}}
    \hfill
    \subfloat{%
        \includegraphics[width=0.31\textwidth]{images/1baseline.png}}
    \hfill
    \subfloat{%
        \includegraphics[width=0.31\textwidth]{images/1lora.png}}

    \vspace{0.5em}

    \subfloat{%
        \includegraphics[width=0.31\textwidth]{images/2final.png}}
    \hfill
    \subfloat{%
        \includegraphics[width=0.31\textwidth]{images/2baseline.png}}
    \hfill
    \subfloat{%
        \includegraphics[width=0.31\textwidth]{images/2lora.png}}

    \vspace{0.5em}

    \subfloat{%
        \includegraphics[width=0.31\textwidth]{images/3final.png}}
    \hfill
    \subfloat{%
        \includegraphics[width=0.31\textwidth]{images/3baseline.png}}
    \hfill
    \subfloat{%
        \includegraphics[width=0.31\textwidth]{images/3lora.png}}

    \vspace{0.5em}

    \subfloat{%
        \includegraphics[width=0.31\textwidth]{images/4final.png}}
    \hfill
    \subfloat{%
        \includegraphics[width=0.31\textwidth]{images/4baseline.png}}
    \hfill
    \subfloat{%
        \includegraphics[width=0.31\textwidth]{images/4lora.png}}

    \vspace{0.5em}

    \makebox[0.31\textwidth][c]{\small Obraz źródłowy}\hfill
    \makebox[0.31\textwidth][c]{\small DDIM Inv.}\hfill
    \makebox[0.31\textwidth][c]{\small LoRA $r=32$}

    \caption{Porównanie rekonstrukcji dla standardowej inwersji DDIM i LoRA $r=32$.}
    \label{fig:qualitative-reconstruction}
\end{figure}

\subsubsection{Aktywacja LoRA w pierwszych krokach}

Po pełnej aktywacji sprawdzono, czy adapter musi działać przez całą
trajektorię inwersji. W tym eksperymencie LoRA była aktywna tylko na
początku trajektorii, a następne kroki wykonywano już zwykłym modelem
bazowym.
Dla każdego rzędu adaptera dostępne są wyniki dla pierwszych pięciu kroków
oraz dla pełnej aktywacji.

\begin{table}[H]
    \centering
    \caption{DDIM Inv., aktywacja LoRA w pierwszych krokach oraz pełna
    aktywacja LoRA dla adapterów $r=8$, $r=16$ oraz $r=32$.}
    \label{tab:lora-rank-window-comparison}
    \resizebox{\textwidth}{!}{%
    \begin{tabular}{lrrrrrrr}
        \toprule
        Adapter &
        Kroki LoRA &
        Corr. $\downarrow$ &
        KL Div. $\times 10^2$ $\downarrow$ &
        Recon. MAE $\downarrow$ &
        LPIPS $\downarrow$ &
        PSNR $\uparrow$ &
        SSIM $\uparrow$ \\
        \midrule
        DDIM Inv. & 0 & 0.1601 & 0.8566 & 0.02663 & 0.09617 & 28.32 & 0.8112 \\
        \midrule
        $r=8$  & 1  & 0.1590 & 0.8499 & 0.02654 & 0.09546 & 28.34 & 0.8123 \\
        $r=8$  & 2  & 0.1204 & 0.3891 & 0.02508 & 0.08507 & 28.65 & 0.8304 \\
        $r=8$  & 3  & 0.1129 & 0.3007 & 0.02489 & 0.08416 & 28.70 & 0.8337 \\
        $r=8$  & 4  & 0.1097 & 0.2582 & 0.02482 & 0.08391 & 28.71 & 0.8349 \\
        $r=8$  & 5  & 0.1078 & 0.2327 & 0.02481 & 0.08393 & 28.71 & 0.8353 \\
        $r=8$  & 50 & 0.0999 & 0.0767 & 0.02063 & 0.05982 & 30.06 & 0.8654 \\
        \midrule
        $r=16$ & 1  & 0.1588 & 0.8499 & 0.02654 & 0.09551 & 28.34 & 0.8122 \\
        $r=16$ & 2  & 0.1180 & 0.3584 & 0.02508 & 0.08475 & 28.64 & 0.8312 \\
        $r=16$ & 3  & 0.1105 & 0.2687 & 0.02486 & 0.08366 & 28.69 & 0.8347 \\
        $r=16$ & 4  & 0.1076 & 0.2302 & 0.02479 & 0.08340 & 28.71 & 0.8359 \\
        $r=16$ & 5  & 0.1059 & 0.2075 & 0.02478 & 0.08346 & 28.70 & 0.8363 \\
        $r=16$ & 50 & \underline{0.0989} & \underline{0.0671} & \underline{0.02028} & \underline{0.05774} & \underline{30.22} & \underline{0.8699} \\
        \midrule
        $r=32$ & 1  & 0.1588 & 0.8486 & 0.02653 & 0.09536 & 28.34 & 0.8124 \\
        $r=32$ & 2  & 0.1143 & 0.3397 & 0.02497 & 0.08402 & 28.67 & 0.8325 \\
        $r=32$ & 3  & 0.1071 & 0.2502 & 0.02474 & 0.08283 & 28.72 & 0.8363 \\
        $r=32$ & 4  & 0.1043 & 0.2096 & 0.02466 & 0.08250 & 28.74 & 0.8377 \\
        $r=32$ & 5  & 0.1028 & 0.1861 & 0.02463 & 0.08249 & 28.74 & 0.8382 \\
        $r=32$ & 50 & \textbf{0.0972} & \textbf{0.0504} & \textbf{0.01970} & \textbf{0.05538} & \textbf{30.40} & \textbf{0.8743} \\
        \bottomrule
    \end{tabular}}
\end{table}

Największy skok jakości pojawia się po przejściu z jednego do dwóch kroków.
Jeden krok prawie nie zmienia wyniku DDIM Inv., natomiast dwa kroki wyraźnie
zmniejszają korelację latentów i jednocześnie poprawiają rekonstrukcję.
Najmocniej widać to dla $r=32$: korelacja spada z $0.1601$ do $0.1143$, KL
Div. $\times 10^2$ z $0.8566$ do $0.3397$, a PSNR rośnie z $28.32$ do
$28.67$. Dalsze zwiększanie liczby kroków z aktywną LoRA daje mniejsze,
ale nadal stabilne zyski.
Wariant $r=32$ aktywny przez pięć kroków osiąga Corr. $0.1028$, KL Div.
$\times 10^2$ równe $0.1861$, PSNR $28.74$ i SSIM $0.8382$.

Pełna aktywacja $r=32$ pozostaje najlepsza liczbowo w tym zestawieniu.
W porównaniu z DDIM Inv. korelacja latentów spada o około $39\%$, KL Div.
$\times 10^2$ o około $94\%$, Recon. MAE o około $26\%$, a LPIPS o około
$42\%$. Równocześnie PSNR wzrasta z $28.32$ do $30.40$, a SSIM z $0.8112$
do $0.8743$. Aktywacja tylko w pierwszych krokach jest jednak istotna
praktycznie: już pierwsze dwa do pięciu kroków przejmują dużą część poprawy
normalności latentów, bez pogorszenia metryk rekonstrukcji.

\subsection{Porównanie z forward diffusion}

Punktem odniesienia dla aktywacji LoRA w pierwszych krokach jest wariant,
w którym pierwsze kroki inwersji zastępowane są zwykłą dyfuzją w przód.
Zestawienie wykonano dla dwóch i pięciu kroków aktywnej LoRA oraz dla
pełnej aktywacji każdego rzędu adaptera.

\begin{table}[H]
    \centering
    \caption{Porównanie aktywacji LoRA w pierwszych krokach z wariantem forward diffusion.}
    \label{tab:forward-reference}
    \resizebox{\textwidth}{!}{%
    \begin{tabular}{lrrrrr}
        \toprule
        Metoda &
        Corr. $\downarrow$ &
        KL Div. $\times 10^2$ $\downarrow$ &
        Recon. MAE $\downarrow$ &
        LPIPS $\downarrow$ &
        PSNR $\uparrow$ \\
        \midrule
        DDIM Inv.             & 0.1601 & 0.8566 & 0.02663 & 0.09617 & 28.32 \\
        Forward 2\%           & 0.1400 & 0.7100 & 0.03000 & 0.10000 & 27.12 \\
        Forward 4\%           & \textbf{0.0900} & 0.3800 & 0.04000 & 0.14000 & 25.68 \\
        \midrule
        LoRA $r=8$, 2 kroki   & 0.1204 & 0.3891 & 0.02508 & 0.08507 & 28.65 \\
        LoRA $r=8$, 5 kroków  & 0.1078 & 0.2327 & 0.02481 & 0.08393 & 28.71 \\
        Pełna LoRA $r=8$      & 0.0999 & 0.0767 & 0.02063 & 0.05982 & 30.06 \\
        \midrule
        LoRA $r=16$, 2 kroki  & 0.1180 & 0.3584 & 0.02508 & 0.08475 & 28.64 \\
        LoRA $r=16$, 5 kroków & 0.1059 & 0.2075 & 0.02478 & 0.08346 & 28.70 \\
        Pełna LoRA $r=16$     & 0.0989 & \underline{0.0671} & \underline{0.02028} & \underline{0.05774} & \underline{30.22} \\
        \midrule
        LoRA $r=32$, 2 kroki  & 0.1143 & 0.3397 & 0.02497 & 0.08402 & 28.67 \\
        LoRA $r=32$, 5 kroków & 0.1028 & 0.1861 & 0.02463 & 0.08249 & 28.74 \\
        Pełna LoRA $r=32$     & \underline{0.0972} & \textbf{0.0504} & \textbf{0.01970} & \textbf{0.05538} & \textbf{30.40} \\
        \bottomrule
    \end{tabular}}
\end{table}

Forward diffusion skutecznie zmniejsza lokalne korelacje w latencie, ale
robi to przez dodanie losowości, co pogarsza rekonstrukcję. Dla wariantu
Forward 4\% korelacja spada do $0.0900$, ale Recon. MAE rośnie do $0.04$,
LPIPS do $0.14$, a PSNR spada do $25.68$. LoRA zachowuje korzystniejszy
kompromis we wszystkich trzech rzędach adaptera: warianty pięciokrokowe
zmniejszają korelację i KL, a jednocześnie utrzymują lepszą rekonstrukcję
niż DDIM Inv. Pełna LoRA $r=32$ osiąga najlepsze metryki rekonstrukcyjne i
najniższą dywergencję KL w tym zestawieniu, natomiast pełna LoRA $r=16$ jest
drugim najlepszym wariantem w metrykach rekonstrukcji.

\subsection{Znaczenie pierwszych kroków inwersji}

Największy wpływ na jakość inwersji mają pierwsze kroki trajektorii.
Wyniki dla wszystkich trzech rzędów adaptera pokazują, że pierwsze dwa
kroki dają największą część poprawy względem DDIM Inv., natomiast aktywacja
adaptera przez pięć pierwszych kroków przejmuje większość zysku pełnej LoRA
w metrykach normalności latentów i edycji tekstowej.

\begin{table}[H]
    \centering
    \caption{Udział poprawy pełnej LoRA uzyskany przez aktywację adaptera
    w pierwszych pięciu krokach. Dla metryk minimalizowanych liczono spadek
    wartości, a dla metryk maksymalizowanych wzrost wartości.}
    \label{tab:first-steps-share}
    \begin{tabular}{lrrrr}
        \toprule
        Metryka &
        $r=8$ &
        $r=16$ &
        $r=32$ &
        Średnia \\
        \midrule
        Corr. & $86.9\%$ & $88.4\%$ & $91.1\%$ & $88.8\%$ \\
        KL Div. $\times 10^2$ & $80.0\%$ & $82.2\%$ & $83.2\%$ & $81.8\%$ \\
        Recon. MAE & $30.4\%$ & $29.2\%$ & $28.8\%$ & $29.5\%$ \\
        LPIPS & $33.7\%$ & $33.1\%$ & $33.6\%$ & $33.4\%$ \\
        CLIP Source & $69.5\%$ & $70.7\%$ & $73.2\%$ & $71.1\%$ \\
        CLIP Target & $81.7\%$ & $81.2\%$ & $82.7\%$ & $81.9\%$ \\
        Directional CLIP & $88.4\%$ & $88.0\%$ & $91.9\%$ & $89.5\%$ \\
        \bottomrule
    \end{tabular}
\end{table}

Wariant pięciokrokowy przejmuje średnio $88.8\%$ poprawy pełnej LoRA dla
korelacji oraz $81.8\%$ poprawy dla dywergencji KL. Podobny efekt widać w
metrykach edycyjnych: średni udział poprawy wynosi $81.9\%$ dla CLIP\cite{DBLP:journals/corr/abs-2103-00020} Target
oraz $89.5\%$ dla Directional CLIP\cite{DBLP:journals/corr/abs-2108-00946}.

Udział poprawy dla metryk rekonstrukcyjnych jest niższy i wynosi średnio
$29.5\%$ dla Recon. MAE oraz $33.4\%$ dla LPIPS. Oznacza to, że pierwsze
kroki są kluczowe dla usuwania nienaturalnej struktury z odtworzonego szumu,
natomiast późniejsze kroki nadal mają znaczenie dla dokładności
rekonstrukcji obrazu. Ten wynik sugeruje praktyczny kompromis jakościowy:
aktywacja LoRA tylko w pierwszych krokach może być wystarczająca, gdy
priorytetem jest edytowalność i normalność latentów, natomiast pełna LoRA
jest korzystniejsza, gdy najważniejsza jest maksymalna jakość rekonstrukcji.


\subsection{Regiony jednorodne i niejednorodne}

Dodatkowo sprawdzono, czy błąd inwersji różni się pomiędzy obszarami
jednorodnymi wizualnie i obszarami zawierającymi więcej lokalnej struktury.
W tym celu obraz podzielono na regiony \emph{plain} i \emph{non-plain}.
Wyniki pokazują, że błąd w regionach jednorodnych jest większy, a odchylenie
standardowe latentów jest tam niższe.

\begin{table}[H]
    \centering
    \caption{Błąd i odchylenie standardowe odtworzonego szumu w regionach
    jednorodnych i niejednorodnych.}
    \label{tab:plain-nonplain}
    \resizebox{\textwidth}{!}{%
    \begin{tabular}{lrrrr}
        \toprule
        Metoda &
        Plain error $\downarrow$ &
        Non-plain error $\downarrow$ &
        Plain std. $\to 1$ &
        Non-plain std. $\to 1$ \\
        \midrule
        DDIM Inv. & 0.1956 & 0.1682 & 0.8770 & 0.9573 \\
        \midrule
        LoRA $r=8$, 1 krok & 0.1937 & 0.1667 & 0.8773 & 0.9576 \\
        LoRA $r=8$, 2 kroki & 0.1305 & 0.1204 & 0.9082 & 0.9750 \\
        LoRA $r=8$, 3 kroki & 0.1196 & 0.1120 & 0.9162 & 0.9801 \\
        LoRA $r=8$, 4 kroki & 0.1148 & 0.1083 & 0.9203 & 0.9830 \\
        LoRA $r=8$, 5 kroków & 0.1121 & 0.1061 & 0.9230 & 0.9849 \\
        Pełna LoRA $r=8$ & 0.1015 & 0.0970 & 0.9433 & \textbf{1.0022} \\
        \midrule
        LoRA $r=16$, 1 krok & 0.1937 & 0.1667 & 0.8773 & 0.9576 \\
        LoRA $r=16$, 2 kroki & 0.1267 & 0.1175 & 0.9104 & 0.9764 \\
        LoRA $r=16$, 3 kroki & 0.1155 & 0.1088 & 0.9188 & 0.9819 \\
        LoRA $r=16$, 4 kroki & 0.1109 & 0.1052 & 0.9228 & 0.9847 \\
        LoRA $r=16$, 5 kroków & 0.1083 & 0.1031 & 0.9252 & 0.9865 \\
        Pełna LoRA $r=16$ & \underline{0.0982} & \underline{0.0944} & \underline{0.9442} & \underline{1.0029} \\
        \midrule
        LoRA $r=32$, 1 krok & 0.1934 & 0.1664 & 0.8774 & 0.9576 \\
        LoRA $r=32$, 2 kroki & 0.1222 & 0.1140 & 0.9122 & 0.9776 \\
        LoRA $r=32$, 3 kroki & 0.1107 & 0.1052 & 0.9208 & 0.9832 \\
        LoRA $r=32$, 4 kroki & 0.1058 & 0.1014 & 0.9251 & 0.9862 \\
        LoRA $r=32$, 5 kroków & 0.1031 & 0.0992 & 0.9278 & 0.9882 \\
        Pełna LoRA $r=32$ & \textbf{0.0930} & \textbf{0.0906} & \textbf{0.9479} & 1.0055 \\
        \bottomrule
    \end{tabular}}
\end{table}

Adapter zmniejsza błąd w obu typach regionów oraz przesuwa odchylenie
standardowe w kierunku wartości $1.0$. Największa poprawa ponownie pojawia
się po przejściu z jednego do dwóch kroków aktywnej LoRA, niezależnie od
rzędu adaptera. Najsilniejszy efekt daje $r=32$, natomiast $r=16$ jest
bardzo bliski w wariancie pełnej aktywacji. Wynik ten jest
spójny z obserwacjami przedstawionymi w pracy \emph{There and Back Again}\cite{staniszewski2026there},
gdzie autorzy wskazują, że problem korelacji i przenoszenia struktury obrazu
do odtworzonego szumu jest szczególnie widoczny w obszarach jednorodnych
wizualnie. W naszym eksperymencie adapter częściowo ogranicza ten efekt,
zmniejszając błąd zarówno w regionach \emph{plain}, jak i \emph{non-plain}.

\subsection{Edycja tekstowa}

Trzeci blok ewaluacji dotyczył przydatności latentów do edycji obrazów. W
tym eksperymencie, po inwersji, obraz był generowany z promptem docelowym.
Prompty docelowe utworzono przez cykliczne przesunięcie listy promptów
źródłowych o jedną pozycję. Taki protokół pozwala porównać warianty inwersji
w tych samych warunkach i sprawdzić, czy uzyskane latenty lepiej reagują na
zmianę warunkowania tekstowego.

\begin{table}[H]
    \centering
    \caption{Metryki CLIP dla edycji tekstowej.}
    \label{tab:clip-editing}
    \resizebox{\textwidth}{!}{%
    \begin{tabular}{lrrr}
        \toprule
        Metoda &
        CLIP Source $\downarrow$ &
        CLIP Target $\uparrow$ &
        Directional CLIP $\uparrow$ \\
        \midrule
        DDIM Inv. & $35.08 \pm 11.73$ & $67.90 \pm 10.92$ & $29.85 \pm 7.46$ \\
        \midrule
        LoRA $r=8$, 1 krok & $35.09 \pm 11.73$ & $67.88 \pm 10.90$ & $29.85 \pm 7.46$ \\
        LoRA $r=8$, 2 kroki & $33.84 \pm 11.77$ & $70.10 \pm 10.38$ & $31.22 \pm 7.22$ \\
        LoRA $r=8$, 3 kroki & $33.45 \pm 11.79$ & $70.74 \pm 10.19$ & $31.53 \pm 7.18$ \\
        LoRA $r=8$, 4 kroki & $33.12 \pm 11.76$ & $71.07 \pm 10.11$ & $31.71 \pm 7.11$ \\
        LoRA $r=8$, 5 kroków & $32.99 \pm 11.72$ & $71.29 \pm 10.05$ & $31.80 \pm 7.10$ \\
        Pełna LoRA $r=8$ & $32.07 \pm 11.82$ & $72.05 \pm 9.72$ & $32.05 \pm 7.01$ \\
        \midrule
        LoRA $r=16$, 1 krok & $35.07 \pm 11.73$ & $67.89 \pm 10.90$ & $29.86 \pm 7.45$ \\
        LoRA $r=16$, 2 kroki & $33.75 \pm 11.80$ & $70.24 \pm 10.37$ & $31.29 \pm 7.20$ \\
        LoRA $r=16$, 3 kroki & $33.35 \pm 11.76$ & $70.82 \pm 10.19$ & $31.59 \pm 7.16$ \\
        LoRA $r=16$, 4 kroki & $33.07 \pm 11.75$ & $71.21 \pm 10.05$ & $31.77 \pm 7.14$ \\
        LoRA $r=16$, 5 kroków & $32.93 \pm 11.74$ & $71.33 \pm 10.06$ & $31.82 \pm 7.11$ \\
        Pełna LoRA $r=16$ & $\underline{32.04 \pm 11.75}$ & $\underline{72.13 \pm 9.71}$ & $\mathbf{32.09 \pm 7.04}$ \\
        \midrule
        LoRA $r=32$, 1 krok & $35.09 \pm 11.74$ & $67.89 \pm 10.89$ & $29.85 \pm 7.45$ \\
        LoRA $r=32$, 2 kroki & $33.75 \pm 11.77$ & $70.24 \pm 10.38$ & $31.29 \pm 7.21$ \\
        LoRA $r=32$, 3 kroki & $33.24 \pm 11.74$ & $70.92 \pm 10.13$ & $31.64 \pm 7.12$ \\
        LoRA $r=32$, 4 kroki & $33.02 \pm 11.72$ & $71.29 \pm 10.08$ & $31.80 \pm 7.09$ \\
        LoRA $r=32$, 5 kroków & $32.83 \pm 11.69$ & $71.42 \pm 10.00$ & $31.89 \pm 7.06$ \\
        Pełna LoRA $r=32$ & $\mathbf{32.00 \pm 11.76}$ & $\mathbf{72.16 \pm 9.76}$ & $\underline{32.07 \pm 7.03}$ \\
        \bottomrule
    \end{tabular}}
\end{table}

Wyniki są spójne z wcześniejszymi obserwacjami.
Warianty LoRA od dwóch kroków wzwyż zwiększają CLIP Target oraz Directional
CLIP, co oznacza silniejsze dopasowanie do kierunku edycji tekstowej. Pełna
LoRA dla wszystkich trzech rzędów osiąga bardzo zbliżone wyniki: najwyższy
CLIP Target uzyskuje $r=32$, a najwyższy Directional CLIP uzyskuje $r=16$.
Jednoczesny spadek CLIP Source wskazuje, że obraz wynikowy w większym
stopniu odchodzi od warunkowania źródłowego. W połączeniu ze wzrostem CLIP
Target i Directional CLIP sugeruje to większą podatność uzyskanych latentów
na zmianę promptu docelowego.


\section{Wyniki dla różnych cfg}

\begin{table}[H]
    \centering
    \caption{Wyniki dla roznych wartosci CFG.}
    \label{tab:cfg-sweep-matched-reconstruction}
    \resizebox{\textwidth}{!}{%
    \begin{tabular}{lrrrrrrrr}
        \toprule
        Metoda &
        CFG &
        Corr. $\downarrow$ &
        KL Div. $\times 10^2$ $\downarrow$ &
        Recon. MAE $\downarrow$ &
        LPIPS $\downarrow$ &
        PSNR $\uparrow$ &
        SSIM $\uparrow$ \\
        \midrule
        DDIM Inv. &
        1.0 & 0.1601 & 0.8566 & 0.02663 & 0.09617 & 28.32 & 0.8112 \\
        LoRA $r=32$ &
        1.0 & 0.0972 & 0.0504 & \textbf{0.01970} & \textbf{0.05538} & \textbf{30.40} & \textbf{0.8743} \\
        \midrule
        LoRA $r=32$ &
        2.0 & 0.0955 & 0.0150 & 0.07249 & 0.29056 & 21.16 & 0.6442 \\
        LoRA $r=32$ &
        3.0 & 0.0989 & 0.0607 & 0.26066 & 0.64225 & 11.36 & 0.3292 \\
        LoRA $r=32$ &
        4.0 & 0.1124 & 0.1500 & 0.32729 & 0.72783 & 9.10 & 0.2355 \\
        LoRA $r=32$ &
        5.0 & 0.1331 & 0.2499 & 0.34274 & 0.75096 & 8.52 & 0.2063 \\
        \bottomrule
    \end{tabular}%
    }
\end{table}


\begin{table}[H]
    \centering
    \caption{Wyniki CFG $=4$ dla danych wygenerowanych przy CFG $=4$.
    (Tylko 100 próbek)}
    \label{tab:cfg4-lora-quick100}
    \resizebox{\textwidth}{!}{%
    \begin{tabular}{llrrrrrrrr}
        \toprule
        Metoda &
        LoRA &
        Probki &
        Std. szumu &
        Corr. $\downarrow$ &
        KL Div. $\downarrow$ &
        Recon. MAE $\downarrow$ &
        LPIPS $\downarrow$ &
        PSNR $\uparrow$ &
        SSIM $\uparrow$ \\
        \midrule
        DDIM Inv. &
        brak & 100 & 0.8812 & 0.3669$^\dagger$ & 2.2262 & 0.1725 & 0.5900 & 13.76 & 0.3566 \\
        DDIM Inv. &
        shared & 100 & \textbf{0.9828} & \textbf{0.3599}$^\dagger$ & \textbf{0.0838} & \textbf{0.0557} & \textbf{0.1939} & \textbf{22.09} & \textbf{0.7522} \\
        DDIM Inv. &
        branch-pair & 100 & 0.9562 & 0.3611$^\dagger$ & 0.3306 & 0.0737 & 0.2489 & 20.40 & 0.7026 \\
        \bottomrule
    \end{tabular}%
    }
\end{table}


\begin{table}[H]
    \centering
    \caption{Przykladowe rekonstrukcje dla CFG $=4$}
    \label{tab:cfg4-qualitative-samples}
    \begin{tabular}{cccc}
        \toprule
        Obraz zrodlowy & DDIM Inv. bez LoRA & Shared LoRA & Branch-pair \\
        \midrule
        \includegraphics[width=0.215\textwidth]{../sample_images_for_cfg/original/final.png} &
        \includegraphics[width=0.215\textwidth]{../sample_images_for_cfg/no_lora/reconstructed.png} &
        \includegraphics[width=0.215\textwidth]{../sample_images_for_cfg/shared_lora/reconstructed.png} &
        \includegraphics[width=0.215\textwidth]{../sample_images_for_cfg/branch/reconstructed.png} \\
        \includegraphics[width=0.215\textwidth]{../sample_images_for_cfg/original/final1.png} &
        \includegraphics[width=0.215\textwidth]{../sample_images_for_cfg/no_lora/reconstructed1.png} &
        \includegraphics[width=0.215\textwidth]{../sample_images_for_cfg/shared_lora/reconstructed1.png} &
        \includegraphics[width=0.215\textwidth]{../sample_images_for_cfg/branch/reconstructed1.png} \\
        \includegraphics[width=0.215\textwidth]{../sample_images_for_cfg/original/final2.png} &
        \includegraphics[width=0.215\textwidth]{../sample_images_for_cfg/no_lora/reconstructed2.png} &
        \includegraphics[width=0.215\textwidth]{../sample_images_for_cfg/shared_lora/reconstructed2.png} &
        \includegraphics[width=0.215\textwidth]{../sample_images_for_cfg/branch/reconstructed2.png} \\
        \includegraphics[width=0.215\textwidth]{../sample_images_for_cfg/original/final3.png} &
        \includegraphics[width=0.215\textwidth]{../sample_images_for_cfg/no_lora/reconstructed3.png} &
        \includegraphics[width=0.215\textwidth]{../sample_images_for_cfg/shared_lora/reconstructed3.png} &
        \includegraphics[width=0.215\textwidth]{../sample_images_for_cfg/branch/reconstructed3.png} \\
        \bottomrule
    \end{tabular}
\end{table}

\begin{table}[H]
    \centering
    \caption{Przykladowe rekonstrukcje dla CFG $=4$}
    \label{tab:cfg4-qualitative-samples-cont}
    \begin{tabular}{cccc}
        \toprule
        Obraz zrodlowy & DDIM Inv. bez LoRA & Shared LoRA & Branch-pair \\
        \midrule
        \includegraphics[width=0.215\textwidth]{../sample_images_for_cfg/original/final4.png} &
        \includegraphics[width=0.215\textwidth]{../sample_images_for_cfg/no_lora/reconstructed4.png} &
        \includegraphics[width=0.215\textwidth]{../sample_images_for_cfg/shared_lora/reconstructed4.png} &
        \includegraphics[width=0.215\textwidth]{../sample_images_for_cfg/branch/reconstructed4.png} \\
        \includegraphics[width=0.215\textwidth]{../sample_images_for_cfg/original/final5.png} &
        \includegraphics[width=0.215\textwidth]{../sample_images_for_cfg/no_lora/reconstructed5.png} &
        \includegraphics[width=0.215\textwidth]{../sample_images_for_cfg/shared_lora/reconstructed5.png} &
        \includegraphics[width=0.215\textwidth]{../sample_images_for_cfg/branch/reconstructed5.png} \\
        \includegraphics[width=0.215\textwidth]{../sample_images_for_cfg/original/final6.png} &
        \includegraphics[width=0.215\textwidth]{../sample_images_for_cfg/no_lora/reconstructed6.png} &
        \includegraphics[width=0.215\textwidth]{../sample_images_for_cfg/shared_lora/reconstructed6.png} &
        \includegraphics[width=0.215\textwidth]{../sample_images_for_cfg/branch/reconstructed6.png} \\
        \bottomrule
    \end{tabular}
\end{table}

\begin{table}[H]
    \centering
    \caption{Porównanie z 4\% (już cały zbiór eval)}
    \label{tab:cfg4-appendix-n3}
    \begin{tabular}{lrrrr}
        \toprule
        Metoda &
        CFG &
        LPIPS $\downarrow$ &
        Corr. $\downarrow$ &
        CLIP/Edit align. $\uparrow$ \\
        \midrule
        DDIM Inv. & 4.0 & 0.525 & 0.172 & 0.725 \\
        Forward $4\%$ & 4.0 & 0.372 & 0.121 & 0.800 \\
        Shared LoRA & 4.0 & \textbf{0.1968} & \textbf{0.0955} & -- \\
        \bottomrule
    \end{tabular}
\end{table}


\section{Divergence w SD 1.5}
\begin{figure}[H]
    \centering
    \includegraphics[width=1\linewidth]{image.png}
    \caption{Val loss }
    \label{fig:placeholder}
\end{figure}


\section{Wnioski}

Przeprowadzone eksperymenty potwierdzają główną hipotezę projektu:
niewielki adapter LoRA używany wyłącznie podczas inwersji może poprawić
jakość odtwarzanego latentu szumowego w DDIM inversion. Adapter zmniejsza
lokalne korelacje w odtworzonym szumie, obniża dywergencję KL względem
rozkładu normalnego oraz zmniejsza błąd względem rzeczywistego szumu
początkowego. Co istotne, poprawa ta nie odbywa się kosztem rekonstrukcji
obrazu.

Najważniejszy wynik dotyczy kompromisu pomiędzy normalnością latentów a
jakością rekonstrukcji. Zarówno wariant aktywny przez pierwsze kilka kroków,
jak i pełna LoRA zmniejszają lokalną strukturę w odtworzonym szumie, a
jednocześnie poprawiają metryki rekonstrukcyjne względem bazowej DDIM
inversion. Oznacza to, że adapter uczy się użytecznej korekty predykcji
szumu, która pozwala uzyskać latent bardziej zbliżony do szumu
gaussowskiego bez widocznego kosztu rekonstrukcyjnego.

Wyniki potwierdzają również szczególne znaczenie początkowych kroków
inwersji. Już aktywacja LoRA przez pierwsze dwa kroki daje największy
pojedynczy przyrost jakości względem DDIM Inv., a pięć pierwszych kroków
przejmuje większość zysku pełnej LoRA w metrykach normalności latentów i
edytowalności. Jest to zgodne z założeniem, że największy błąd aproksymacji
DDIM inversion pojawia się na początku trajektorii.

Porównanie z forward diffusion pokazuje, że poprawa uzyskana przez LoRA ma
inny charakter niż proste dodanie losowego szumu. Forward diffusion
skutecznie obniża korelacje w latencie, ale pogarsza rekonstrukcję. LoRA
natomiast poprawia własności latentów i jednocześnie poprawia metryki
rekonstrukcyjne. Sugeruje to, że adapter uczy się użytecznej korekty
predykcji szumu, a nie tylko zwiększa losowość odtwarzanego latentu.

W eksperymencie edycji tekstowej warianty LoRA od dwóch kroków wzwyż
poprawiają CLIP Target oraz Directional CLIP względem bazowej DDIM inversion.
Oznacza to, że uzyskane latenty są bardziej podatne na zmianę promptu
docelowego.

Podsumowując, wariant aktywny przez pierwsze kilka kroków stanowi
interesujący kompromis pomiędzy normalnością latentów, edytowalnością i
jakością rekonstrukcji, szczególnie dla adaptera $r=32$. W porównaniu
rzędów adaptera najlepszą pełną LoRA był wariant $r=32$, który uzyskał
najniższą korelację, najniższą dywergencję KL oraz najlepsze metryki
rekonstrukcyjne. Projekt pokazuje, że uczenie korekty predykcji szumu jest
sensownym kierunkiem poprawy DDIM inversion i może prowadzić do latentów
lepiej nadających się zarówno do rekonstrukcji, jak i do dalszej edycji
tekstowej.

\newpage
\printbibliography
\end{document}
